\scanpage{351}
\uline{Non-Abelian} \uline{periods} \uline{on} \uline{smooth} \uline{manifolds}

We can find from the $d$-formulae of the model a set
of integrals defining a determining set of periods for
elements of the nilpotent $\pi_1$ quotients $\{\Gamma_k\}$. There is an
analogous discussion for the higher homotopy groups.

For example suppose the model begins as $\Lambda(x_1,x_2,y)$
with $dy=x_1\wedge x_2$. Let $\omega_1,\omega_2,\eta$ denote the 1-forms in
the manifold which are the $\rho$-images of these model generators.
If $\gamma$ is a loop in $M$ based at $p$ we can form the three
ingegrals
\[
\int_\gamma\omega_1,\ \int_\gamma\omega_2,\
\int_\gamma\eta-\int_\gamma\left(\int_{\gamma_t}\omega_1\right)\omega_2
\quad{}^1
\]

The first two periods are homology invariants of
since $\omega_1$ and $\omega_2$ are closed. The third integral is a
homotopy invariant of $\gamma$ because $d\eta=\omega_1\wedge\omega_2$. If we
traverse $\gamma_1$ and then $\gamma_2$ the periods multiply like the
matrices
\[
\begin{pmatrix}1&x&z\\0&1&\unsure{\uline{y}|y}\\0&0&1\end{pmatrix}.
\]

\figure{\begin{tikzpicture}\draw (0,0)--(5,0);\end{tikzpicture}}

1\quad These iterated integrals have been considered in
a general way, by Chen. Our theory seems to guide
their efficient application to homotopy theory.

\mspage{352}
